\chapter{Controller Source Listings}
\label{app:code}

This appendix collects the source listings of the implemented Cartesian
impedance controller. The code is taken from the controller used for the
table-contact experiment and is referenced from
Chapter~\ref{ch:implementation} where the corresponding equations are derived.
The listings use the C++ types \texttt{Vec3}, \texttt{Vec6}, \texttt{Vec7},
\texttt{Mat3}, and \texttt{Mat6x7} as aliases for the corresponding fixed-size
Eigen matrices, and \texttt{Parameters} for the controller configuration
struct.

\section{Surface Task Frame and Spatial Gains}
\label{app:surface_gains}

The surface task frame is built from a plane normal and a tangent. The spatial
gain matrix expresses diagonal surface-frame gains in the base frame, as in
\cref{eq:spatial_gain_matrix}.

\begin{lstlisting}[style=cppstyle, language=C++, caption={Surface task frame and spatial gain matrix.}, label={lst:app_surface_gains}]
inline Mat3 makeSurfaceFrameFromNormalTangent(
    const Vec3& normal_input, const Vec3& tangent1_input) {
  const Vec3 normal = normalizedOrFallback(normal_input, Vec3(1.0, 0.0, 0.0));
  Vec3 tangent1 = tangent1_input - normal * normal.dot(tangent1_input);

  if (tangent1.norm() <= 1e-9) {
    Vec3 fallback(0.0, 1.0, 0.0);
    if (std::abs(normal.dot(fallback)) > 0.95) {
      fallback = Vec3(0.0, 0.0, 1.0);
    }
    tangent1 = fallback - normal * normal.dot(fallback);
  }

  tangent1.normalize();
  Vec3 tangent2 = normal.cross(tangent1);
  tangent2.normalize();

  Mat3 R_surface;
  R_surface.col(0) = normal;
  R_surface.col(1) = tangent1;
  R_surface.col(2) = tangent2;
  return R_surface;
}

inline Mat3 makeSpatialGainMatrix(
    const Vec3& diagonal_in_surface_frame, const Mat3& R_surface) {
  return R_surface * diagonal_in_surface_frame.asDiagonal()
         * R_surface.transpose();
}
\end{lstlisting}

\section{Plane Projection of the Desired Motion}
\label{app:desired_motion}

The desired position follows a smooth trajectory in free space and is projected
onto the surface plane when the constraint is enabled, as in
\cref{eq:signed_normal_distance,eq:plane_projected_desired_position}.

\begin{lstlisting}[style=cppstyle, language=C++, caption={Desired motion with plane projection.}, label={lst:app_desired_motion}]
inline DesiredMotion computeDesiredMotion(
    const Parameters& params,
    double time,
    const Vec3& p_start,
    const Vec3& p_EE,
    const Mat3& R_surface,
    const Vec3& plane_point) {
  DesiredMotion desired{p_start, Vec3::Zero()};

  if (!params.hold_mode) {
    const double r = time / params.trajectory_duration;
    const double s = smoothStep(r);
    const double s_dot = smoothStepDerivative(r, params.trajectory_duration);
    desired.p_d = p_start + s * params.delta_p;
    desired.pdot_d = s_dot * params.delta_p;
  }

  if (params.constraint_enabled) {
    const Vec3 normal = R_surface.col(0);
    // signed_distance = n^T (p_plane - p_EE); p_d = p_EE + signed_distance * n
    const double signed_distance = normal.dot(plane_point - p_EE);
    desired.p_d = p_EE + signed_distance * normal;
    desired.pdot_d.setZero();
  }

  return desired;
}
\end{lstlisting}

\section{Orientation Error and Rotational Constraint}
\label{app:orientation}

The orientation error is returned in the base frame
(\cref{eq:rotation_error_vector}), and the rotational constraint is applied in
the surface frame (\cref{eq:surface_rotational_mask}).

\begin{lstlisting}[style=cppstyle, language=C++, caption={Base-frame orientation error and surface-frame rotational mask.}, label={lst:app_orientation}]
inline Vec3 orientationError(const Mat3& R_current, const Mat3& R_desired) {
  Mat3 R_error = R_current.transpose() * R_desired;
  Eigen::AngleAxisd angle_axis(R_error);

  if (std::abs(angle_axis.angle()) < 1e-9) {
    return Vec3::Zero();
  }
  // Rotate the axis-angle error into the base frame.
  return R_current * angle_axis.axis() * angle_axis.angle();
}

inline Vec3 applyRotationalAxisMask(
    const Parameters& params, Vec3 e_R, const Mat3& R_surface) {
  if (!params.constraint_enabled) {
    return e_R;
  }

  Vec3 e_R_task = R_surface.transpose() * e_R;
  e_R_task(0) = params.constrain_rotation_about_surface_normal   ? e_R_task(0) : 0.0;
  e_R_task(1) = params.constrain_rotation_about_surface_tangent1 ? e_R_task(1) : 0.0;
  e_R_task(2) = params.constrain_rotation_about_surface_tangent2 ? e_R_task(2) : 0.0;

  return R_surface * e_R_task;
}
\end{lstlisting}

\section{Null-Space Torque}
\label{app:nullspace}

The null-space torque uses the damped projector of the full Jacobian, a
return-to-start spring, joint-space damping, and the smallest-singular-value
optimization, as in \cref{eq:implementation_full_projector,eq:implementation_nullspace_torque}.

\begin{lstlisting}[style=cppstyle, language=C++, caption={Full-Jacobian null-space torque with return-to-start and sigma-min optimization.}, label={lst:app_nullspace}]
inline double smallestSingularValue(const Mat6x7& J) {
  Eigen::JacobiSVD<Mat6x7> svd(J, Eigen::ComputeFullU | Eigen::ComputeFullV);
  return svd.singularValues().minCoeff();
}

inline Vec7 limitJointTorqueVectorNorm(const Vec7& v, double max_norm) {
  if (max_norm <= 0.0) {
    return v;
  }
  const double norm = v.norm();
  if (norm > max_norm && norm > 1e-12) {
    return max_norm * v / norm;
  }
  return v;
}

inline Vec7 computeNullspaceTorque(
    const Parameters& params,
    const Model& model,
    const RobotState& state,
    const Mat6x7& J,
    const Vec7& dq,
    const Vec7& q_start) {
  Vec7 tau_nullspace = Vec7::Zero();
  if (!params.use_nullspace_optimization) {
    return tau_nullspace;
  }

  Map<const Vec7> q_current(state.q.data());
  const Mat7x7 I7 = Mat7x7::Identity();
  const Mat6x6 I6 = Mat6x6::Identity();
  const double lambda = 0.05;

  // Damped null-space projector of the full Jacobian:
  //   N = I - J^T (J J^T + lambda^2 I)^(-1) J
  const Mat6x6 JJt_damped = J * J.transpose() + lambda * lambda * I6;
  const Mat7x7 N = I7 - J.transpose() * JJt_damped.ldlt().solve(J);

  // Return-to-start spring and joint-space damping, projected into the null space.
  tau_nullspace = N * (params.nullspace_k_start * (q_start - q_current)
                       - params.nullspace_damping * dq);

  // Smallest-singular-value direction: last right singular vector of J.
  Eigen::JacobiSVD<Mat6x7> svd_current(J, Eigen::ComputeFullU | Eigen::ComputeFullV);
  Vec7 n = svd_current.matrixV().col(6);
  if (n.norm() <= 1e-9) {
    return limitJointTorqueVectorNorm(tau_nullspace, params.nullspace_tau_max);
  }
  n.normalize();

  // Compare sigma_min at the two candidate configurations q +/- alpha n.
  const Vec7 q_plus = q_current + params.nullspace_alpha * n;
  const Vec7 q_minus = q_current - params.nullspace_alpha * n;
  const std::array<double, 42> J_plus_array = model.zeroJacobian(
      Frame::kEndEffector, vec7ToArray(q_plus), state.F_T_EE, state.EE_T_K);
  const std::array<double, 42> J_minus_array = model.zeroJacobian(
      Frame::kEndEffector, vec7ToArray(q_minus), state.F_T_EE, state.EE_T_K);
  Map<const Mat6x7> J_plus(J_plus_array.data());
  Map<const Mat6x7> J_minus(J_minus_array.data());

  const double sigma_min_plus = smallestSingularValue(J_plus);
  const double sigma_min_minus = smallestSingularValue(J_minus);
  const double sigma_direction = (sigma_min_plus > sigma_min_minus) ? 1.0 : -1.0;

  const Vec7 tau_sigma = N * (params.nullspace_k_sigma * sigma_direction
                             * params.nullspace_alpha * n);

  return limitJointTorqueVectorNorm(tau_nullspace + tau_sigma,
                                    params.nullspace_tau_max);
}
\end{lstlisting}

\section{Core Control Callback}
\label{app:callback_core}

The core of the real-time callback evaluates the impedance wrench, maps it to
joint torques, adds the projected null-space torque and the Coriolis term, and
returns the command. The gain matrices \texttt{Kp\_used}, \texttt{KR\_used},
etc. are the phase-dependent surface-frame gains selected by the control phase.

\begin{lstlisting}[style=cppstyle, language=C++, caption={Core of the real-time impedance-control callback.}, label={lst:app_callback_core}]
// Differential kinematics and current pose.
Vec6 xdot = J * dq;
Vec3 pdot = xdot.head<3>();
Vec3 omega = xdot.tail<3>();
Vec3 p_EE = T_EE.block<3, 1>(0, 3);
Mat3 R_EE = T_EE.block<3, 3>(0, 0);

// Desired motion: smooth trajectory in free space, plane projection in contact.
DesiredMotion desired = computeDesiredMotion(
    params, time, p_start, p_EE, R_surface, plane_point);
Vec3 e_p = desired.p_d - p_EE;

// Base-frame orientation error, masked in the surface frame.
Vec3 e_R = applyRotationalAxisMask(
    params, orientationError(R_EE, R_d_phase), R_surface);

// Impedance wrench with phase-dependent surface-frame gains.
Vec3 f = Kp_used * e_p + Dp_used * (desired.pdot_d - pdot);
Vec3 m = KR_used * e_R - DR_used * omega;
Vec6 wrench;
wrench.head<3>() = f;
wrench.tail<3>() = m;

Vec7 tau_task = J.transpose() * wrench;

// Null-space torque is enabled once contact search and orientation are done.
Vec7 tau_nullspace = Vec7::Zero();
if (nullspace_active) {
  tau_nullspace = computeNullspaceTorque(params, model, state, J, dq, q_start);
}

Array7 coriolis_array = model.coriolis(state);
Map<const Vec7> coriolis(coriolis_array.data());

Vec7 tau_cmd = tau_task + tau_nullspace + coriolis;
\end{lstlisting}
